% Generated draft: requires mathematical and source review.
\documentclass[11pt]{article}
\usepackage[margin=25mm]{geometry}
\usepackage{fontspec}
\setmainfont{Latin Modern Roman}
\usepackage{amsmath,amssymb,mathtools,unicode-math}
\setmathfont{Latin Modern Math}
\usepackage{enumitem,needspace,xurl,hyperref}
\setlength{\parindent}{0pt}
\setlength{\parskip}{5pt}
\setlength{\emergencystretch}{3em}
\newcommand{\N}{\mathbb{N}}
\newcommand{\Z}{\mathbb{Z}}
\newcommand{\Q}{\mathbb{Q}}
\newcommand{\R}{\mathbb{R}}
\newcommand{\C}{\mathbb{C}}
\newcommand{\F}{\mathbb{F}}
\newcommand{\A}{\mathbb{A}}
\newcommand{\PP}{\mathbb{P}}
\newcommand{\E}{\mathbb{E}}
\newcommand{\eps}{\varepsilon}
\newcommand{\dd}{\,\mathrm{d}}
\begin{document}
\section*{UnsolvedMath ID: 700154}
\textbf{Source problem code:} catalogue entry
\subsection{Definitions and mathematical statement}
For subgroups \(X,Y\) of a group, let \([X,Y]\) denote the subgroup generated by the commutators \(x^{-1}y^{-1}xy\), where \(x\in X\) and \(y\in Y\). A group \(H\) is nilpotent if its lower central series, defined by \(\gamma_1(H)=H\) and \(\gamma_{n+1}(H)=[\gamma_n(H),H]\), satisfies \(\gamma_{c+1}(H)=\{1\}\) for some integer \(c\geq 1\). It is soluble if its derived series, defined by \(H^{(0)}=H\) and \(H^{(n+1)}=[H^{(n)},H^{(n)}]\), satisfies \(H^{(d)}=\{1\}\) for some integer \(d\geq 0\).

A group is locally finite, locally nilpotent, or locally soluble if every finitely generated subgroup is, respectively, finite, nilpotent, or soluble. The nilpotency class or derived length may depend on the finitely generated subgroup.

The question proposed by V. P. Shunkov is whether the following assertion holds: for every locally finite group \(G\) and all locally nilpotent subgroups \(A,B\leq G\) such that
\[
G=AB=\{ab:a\in A,\ b\in B\},
\]
the group \(G\) is locally soluble. Equivalently, must \(\langle S\rangle\) be soluble for every finite subset \(S\subseteq G\)? No normality or commutativity assumption on the factors is imposed in the supplied statement.
\subsection{Short English statement}
Must every locally finite group that is the setwise product of two locally nilpotent subgroups be locally soluble? This is the supplied question of V. P. Shunkov, recorded as Kourovka Notebook Problem 2.84.
\subsection{Sources}
\begin{itemize}
\item Supplied local UnsolvedMath registry record: permanent ID \(700154\), exact source code \texttt{catalogue entry}, title \emph{Kourovka Notebook Problem 2.84}; upstream revision \texttt{d371ee917ac9647fbb06cd6fc529deec6f3dda82}.
\item The supplied record attributes the problem to V. P. Shunkov, proposed in 1967, and cites \emph{The Kourovka Notebook (Unsolved Problems in Group Theory)}, 2nd issue (1967), problem 2.84, as recorded in the 21st-issue edition (2026), \texttt{arXiv:1401.0300v48}. The supplied collection description names E. I. Khukhro and V. D. Mazurov as editors.
\item The registry labels the problem \texttt{open} and records an update timestamp of \texttt{2026-10-16T14:00:00Z}. This is supplied registry metadata, not an independently verified assessment of the literature or mathematical status. The cited source has not been independently consulted for this draft.
\end{itemize}
\end{document}
